\section{Conclusion} \label{sec:conclusion_and_discussion}  This work revisits the common assumption that GRPO must remain in a low-staleness, near-on-policy regime to train LLMs effectively. We show that this constraint is overly conservative: stale rollouts can retain useful learning signal far beyond the standard regime, but high-staleness optimization fails unless its localized instability is handled explicitly. Our diagnosis identifies negative-advantage suffix updates after an off-support boundary as a key failure mode, motivating $\mu$-GRPO, a staged framework that decouples rollout generation from policy optimization and stabilizes large-$\mu$ training with relaxed clipping and negative-advantage veto. Across multiple models and mathematical reasoning benchmarks, $\mu$-GRPO converts GRPO's staleness tolerance into a substantially improved performance--efficiency trade-off, achieving up to a 2.2$\times$ wall-clock speedup while preserving or improving final performance. Our study focuses primarily on math RLVR, with additional coding and asynchronous evaluations. While exact wall-clock gains depend on hardware and implementation, the reduction in rollout refreshes and synchronization follows directly from the large-stage design. Extending high-staleness training to broader domains and larger models remains future work.